% Korean cover text and fonts, read by book/cover.tex (`make cover`, ...).
% The layout lives in book/cover.tex; only words and fonts go here.
\usepackage[hangul]{xetexko}
\setmainfont{Noto Serif CJK KR}\setmainhangulfont{Noto Serif CJK KR}
\setsansfont{Noto Sans CJK KR}\setsanshangulfont{Noto Sans CJK KR}

\newcommand\booktitle{프로덕트를 처음부터 디자인하는 방법}
\newcommand\bookauthor{박근우}

% Front cover: the title as broken into lines, and the cycle diagram.
\newcommand\fronttitle{프로덕트를\\처음부터\\디자인하는 방법}
\newcommand\cycledesign{디자인}
\newcommand\cyclebuild{구현}
\newcommand\cycleevaluate{평가}
\newcommand\cyclecenter{끝없는\\반복}

% Back cover.
\newcommand\backtext{%
        이 소책자는 사용자 중심으로 제품을 디자인하는 방법을 설명한다.
        상황 파악에서 시작해 니즈 도출, 스케치, 프로토타이핑,
        그리고 디자인 커뮤니케이션과 리포트 작성까지,
        챕터를 따라가면 제품 하나를 디자인할 수 있도록 구성했다.
        마지막으로 AI로 인해 변하는 것과 변하지 않는 것을 이야기한다.\par\medskip
        \textsf{\bfseries 이런 분들에게}\par\smallskip
        · UX 디자이너로 일을 하고 싶은 학생 또는 디자이너\par
        · 구현뿐만 아니라 서비스 설계까지 하고 싶은 개발자\par
        · 사용자 중심으로 제품을 만들고 싶은 모든 사람}
